\documentclass[11pt]{amsart}
\usepackage[margin=1.1in]{geometry}
\usepackage{amsmath,amssymb,amsthm}
\usepackage{booktabs}
\usepackage{xcolor}
\usepackage{hyperref}
\hypersetup{colorlinks=true,linkcolor=blue!50!black,citecolor=blue!50!black,urlcolor=blue!50!black,pdftitle={The density of admissible n in Erdos-Selfridge problem 1063},pdfauthor={Paul Lee}}

\newtheorem{theorem}{Theorem}[section]
\newtheorem{lemma}[theorem]{Lemma}
\newtheorem{corollary}[theorem]{Corollary}
\newtheorem{proposition}[theorem]{Proposition}
\newtheorem{conjecture}[theorem]{Conjecture}
\theoremstyle{remark}
\newtheorem{remark}[theorem]{Remark}

\DeclareMathOperator{\Bad}{Bad}
\DeclareMathOperator{\rad}{rad}
\renewcommand{\Pr}{\operatorname{Pr}}
\newcommand{\Exp}{\mathrm{Exp}}
\newcommand{\N}{\mathbb N}

\title[The density of admissible $n$ in Erd\H{o}s problem 1063]{The density of admissible $n$ in Erd\H{o}s--Selfridge problem 1063, with a growth conjecture for $n_k$}
\author{Paul Lee}
\address{Houston, Texas}
\date{October 23, 2026}

\begin{document}

\begin{abstract}
For $k\ge2$ let $n_k\ge2k$ be the least $n$ such that $n-i$ divides $\binom nk$ for all but exactly one $0\le i<k$ (Erd\H{o}s problem 1063, OEIS A389360). We describe, prime by prime, which index fails: $(n-i)\nmid\binom nk$ iff $v_p(n-i)>M_p(n)-d_p(n)$ for some prime $p\le k$, where $M_p$ is the largest $p$-adic valuation in the block $\{n-k+1,\dots,n\}$ and $d_p$ counts the non-carries in the lowest $\lfloor\log_pk\rfloor$ digits when $k$ and $n-k$ are added in base $p$. This gives a two-line proof of the Erd\H{o}s--Selfridge observation and the explicit lower bound $n_k-i_0\ge\prod_{p^a\|k}p^{a+\lfloor\log_p(k-1)\rfloor}$ for the failing index $i_0$ (with equality for $k=4,6,14$), hence $n_k\ge\prod_{p^a\|k}p^{a+\lfloor\log_p(k-1)\rfloor}$; this bound is formally verified in Lean~4. The admissible $n$ form a periodic set whose density $\delta(k)$ is an explicit finite product-sum, and we prove
\[\log\frac1{\delta(k)}=(C+o(1))\,\frac{k}{\log k},\qquad C=\sum_{j\ge1}\frac1j\log\Big(1+\frac1j\Big)=1.25774688694\ldots\]
We conjecture that $n_k\delta(k)$ behaves like an exponential random variable, so that $\log n_k=(C+o(1))k/\log k$, and support this with the values $n_k$ for $60\le k\le140$ and for prime $k\le167$ (extending the OEIS table, which ends at $k=59$, and the values for $k\le75$ posted in the discussion thread), four pre-registered out-of-sample tests, and a correlation of $0.993$ between $\log n_k$ and $\log(1/\delta(k))$ for $31\le k\le140$. The conjecture lies strictly between the bounds $c(\log k)^2$ and $(1+o(1))\,k\log\log k/\log k$ recently claimed by Cipollini, and predicts that the upper one is off by a factor of order $\log\log k$.
\end{abstract}

\maketitle

\section{Introduction}

Erd\H{o}s and Selfridge \cite{ErSe83} observed that for $n\ge2k$ at least one of the $k$ consecutive integers $n,n-1,\dots,n-k+1$ fails to divide $\binom nk$; a proof was given by Monier \cite{Mo85}. Problem 1063 on \texttt{erdosproblems.com} \cite{EP1063} (also B31 in Guy's collection \cite{Gu04}) asks to estimate
\[n_k=\min\{\,n\ge2k:\ (n-i)\mid\tbinom nk\text{ for all but exactly one }0\le i<k\,\}.\]
Monier showed $n_k\le k!$ for $k\ge3$, and Cambie improved this to $n_k\le k\cdot\operatorname{lcm}(2,\dots,k-1)=e^{(1+o(1))k}$. The OEIS entry A389360 \cite{OEIS} lists $n_k$ for $k\le59$. In the problem's discussion thread \cite{EP1063}, the user ``rickyc'' proved $n_k\ge\max\big(2k,\prod_{p^a\|k}p^{a+\lfloor\log_p(k/2)\rfloor}\big)$ (26 June 2026, with GPT-5.5 Thinking), and then the $k-1$ version in the sharper form $p^{\,a+v_p\binom{k-1}{e}+\lfloor\log_p\max(e,k-1-e)\rfloor}\mid n-e$ for the failing index $e$ (10 July 2026, with GPT-5.6 Sol). Our Proposition~\ref{prop:5} is an independent proof of the $k-1$ bound; what is new here is the identification of the equality cases, the characterisation of the bad set behind it, and a formal verification in Lean~4. Two proof claims by Cipollini \cite{Ci26a,Ci26b}, listed on the problem page (which notes that listing is no guarantee of correctness), assert
\begin{equation}\label{eq:cip}
c(\log k)^2\ \le\ \log n_k\ \le\ \frac{k}{\log k}\Big(\log\log k+\log\log\log k+\log2+o(1)\Big).
\end{equation}
A bound of the same shape, $\log n_k\le Ck\log\log k/\log k$, was announced independently two days earlier by White \cite{Wh26}; an upper bound of exactly the form $n_k\ll\exp\big(\frac{k}{\log k}(\log\log k+\log\log\log k+\log 2)\big)$ is recorded in DeepMind's formal-conjectures repository as formally proved in Lean~4 by the LEAP agent \cite{Ku26,FC}; we have not re-checked that proof. No lower bound for $n_k$ appears among the formalized results there.

Throughout, $k\ge2$ and $n\ge2k$ are integers, the \emph{block} is $B=\{n-k+1,\dots,n\}$, and the index $i\in[0,k)$ refers to the element $n-i$. We call $i$ \emph{bad} if $(n-i)\nmid\binom nk$, and $n$ \emph{admissible} if it has exactly one bad index. For a prime $p$ and an integer $x\ne0$, $v_p(x)$ is the $p$-adic valuation, and
\[L_p=L_p(k)=\lfloor\log_pk\rfloor,\qquad M_p(n)=\max_{0\le i<k}v_p(n-i),\]
\[d_p(n)=\#\{1\le e\le L_p:\ n\bmod p^e\ \ge\ k\bmod p^e\}.\]
Note that $d_p(n)$ counts the levels $e\le L_p$ at which adding $k$ and $n-k$ in base $p$ produces no carry.

Our results are the following.

\begin{theorem}[local description]\label{thm:local}
For every prime $p\le k$, $v_p\binom nk=M_p(n)-d_p(n)$, and an index $i$ is bad if and only if $v_p(n-i)>M_p(n)-d_p(n)$ for some prime $p\le k$. If $p\mid k$ then $d_p(n)\ge v_p(k)\ge1$, so every index of largest $p$-adic valuation is bad; in particular some index is always bad.
\end{theorem}

\begin{theorem}[explicit lower bound]\label{thm:lower}
If $n$ is admissible with bad index $i_0$, then
\[W'_k:=\prod_{p^a\|k}p^{\,a+\lfloor\log_p(k-1)\rfloor}\quad\text{divides}\quad n-i_0.\]
Hence $n_k\ge W'_k$. Equality $n_k-i_0=W'_k$ holds for $k=4,6,14$ and for no other $k\le140$.
\end{theorem}

\begin{theorem}[exact density]\label{thm:density}
Admissibility of $n$ depends only on $n\bmod p^{2L_p}$ for the primes $p\le k$. The admissible integers therefore form a union of residue classes modulo $D_k=\prod_{p\le k}p^{2L_p}$, of density
\begin{equation}\label{eq:delta}
\delta(k)=\sum_{i_0=0}^{k-1}\ \prod_{p\le k}\Pr\big(E_{p,i_0}\big),\qquad E_{p,i_0}=\{n:\ \Bad_p(n)\subseteq\{i_0\}\},
\end{equation}
where $\Bad_p(n)=\{i:\ v_p(n-i)>M_p(n)-d_p(n)\}$ and $\Pr(E_{p,i_0})$ is the proportion of residue classes modulo $p^{2L_p}$ in $E_{p,i_0}$. For $p>\sqrt k$,
\begin{equation}\label{eq:closed}
\Pr(E_{p,i_0})=\frac{k\bmod p}{p}+\big[\,i_0\bmod p\ \ge\ k\bmod p\,\big]\cdot\begin{cases}1/p&\text{if }\lfloor k/p\rfloor=1,\\ 1/p^2&\text{if }\lfloor k/p\rfloor\ge2.\end{cases}
\end{equation}
\end{theorem}

\begin{theorem}[size of the density]\label{thm:asymp}
Let $S(k)=\sum_{\sqrt k<p\le k,\ p\nmid k}\log\dfrac{p}{k\bmod p}$. Then
\[\log\frac1{\delta(k)}=S(k)+O(\sqrt k)\qquad\text{and}\qquad S(k)=(C+o(1))\,\frac{k}{\log k},\]
where
\[C=\int_1^\infty\frac{-\log\{x\}}{x^2}\,dx=\sum_{j\ge1}\frac1j\log\Big(1+\frac1j\Big)=\sum_{m\ge1}\frac{(-1)^{m+1}\zeta(m+1)}{m}=1.25774688694436963\ldots\]
Consequently $\log(1/\delta(k))=(C+o(1))\,k/\log k$.
\end{theorem}

The density is a statement about the set of admissible $n$, not about its least element. The following is a conjecture.

\begin{conjecture}\label{conj}
$n_k\,\delta(k)$ is approximately an $\Exp(1)$ random variable, so that typically $\log n_k=\log(1/\delta(k))+O(1)$. Under this heuristic, and only under it,
\[\log n_k=(C+o(1))\,\frac{k}{\log k}.\]
\end{conjecture}

Section~\ref{sec:evidence} reports the evidence: the values $n_k$ for $60\le k\le140$ and for prime $k\le167$ (Tables~\ref{tab:values} and~\ref{tab:values2}), four pre-registered out-of-sample tests in which $78$ of $86$ values fell inside the central $90\%$ interval of $\Exp(1)/\delta(k)$ ($77.4$ expected), and the statistics of $\log n_k-\log(1/\delta(k))$. Section~\ref{sec:cip} explains why the construction behind the upper bound in \eqref{eq:cip} loses a factor of order $\log\log k$ relative to the conjecture.

\subsection*{Provenance}
The analysis, the search code and the first draft of this paper were produced with Claude (Anthropic) at the author's direction. The mathematics was independently refereed by further AI reviewers, including Claude Fable 5.1, Gemini 4 Argon and a Meta AI model, using independently written code; every value in Tables~\ref{tab:values} and~\ref{tab:values2} was verified by exact integer arithmetic; minimality was confirmed by two independently written exhaustive searches for every $k\le119$ and for $k=127$, and for the remaining values rests on one search that reproduces every previously computed value (Section~\ref{sec:evidence}). Code, data and the pre-registered prediction files are available at \cite{code}.

\section{The bad set, prime by prime}\label{sec:local}

\begin{lemma}\label{lem:1}
For $0\le i<k$: $(n-i)\mid\binom nk$ iff $v_p(n-i)\le v_p\binom nk$ for every prime $p\le k$.
\end{lemma}
\begin{proof}
With $P=\prod_{j<k}(n-j)=k!\binom nk$ we have $(n-i)\mid\binom nk\iff k!\mid P/(n-i)$. For $p>k$, $v_p(k!)=0$, so there is no condition. For $p\le k$ the condition $v_p(P)-v_p(n-i)\ge v_p(k!)$ reads $v_p(n-i)\le v_p(P)-v_p(k!)=v_p\binom nk$.
\end{proof}

So the set of bad indices is $\bigcup_{p\le k}\Bad_p(n)$ with $\Bad_p(n)=\{i: v_p(n-i)>v_p\binom nk\}$.

\begin{lemma}\label{lem:2}
For $q=p^e$, the block contains $\lfloor k/q\rfloor+\delta_e$ multiples of $q$, where $\delta_e=\delta_e(n)=[\,n\bmod q<k\bmod q\,]$. Hence $v_p\binom nk=\sum_{e\ge1}\delta_e$.
\end{lemma}
\begin{proof}
Write $n=aq+r$ and $k=bq+s$ with $0\le r,s<q$. The number of multiples of $q$ in the block is $\lfloor n/q\rfloor-\lfloor(n-k)/q\rfloor$, which equals $b$ if $r\ge s$ and $b+1$ if $r<s$. Summing over $e$ and using Legendre's formula, $v_p(P)=\sum_e(\lfloor k/p^e\rfloor+\delta_e)=v_p(k!)+\sum_e\delta_e$.
\end{proof}

Thus $d_p(n)=\#\{1\le e\le L_p:\ \delta_e(n)=0\}$.

\begin{lemma}\label{lem:3}
$v_p\binom nk=M_p(n)-d_p(n)$. Consequently:
\begin{enumerate}
\item $\Bad_p(n)=\varnothing$ iff $d_p(n)=0$;
\item if $d_p(n)\ge1$, every index of maximal $p$-adic valuation is bad, and $|\Bad_p(n)|=1$ iff the maximum is attained at a single index $i^*$ and every other index has valuation $\le M_p-d_p$.
\end{enumerate}
\end{lemma}
\begin{proof}
Write $L=L_p$. The block has $k\ge p^L$ elements, so it contains a multiple of $p^L$ and $M_p\ge L$. For $e>L$ we have $p^e>k$, so $\delta_e=[\text{the block contains a multiple of }p^e]=[M_p\ge e]$. Hence $\sum_e\delta_e=(L-d_p)+(M_p-L)=M_p-d_p$. Parts (1) and (2) follow from $\Bad_p(n)=\{i: v_p(n-i)>M_p-d_p\}$.
\end{proof}

\begin{corollary}[Erd\H{o}s--Selfridge]\label{cor:4}
If $p\mid k$ then $d_p(n)\ge v_p(k)\ge1$ for every $n$, because $k\bmod p^e=0$ for $e\le v_p(k)$ (and $v_p(k)\le L_p$). So $\Bad_p(n)\ne\varnothing$, and for every $n\ge k$ some $n-i$ with $0\le i<k$ fails to divide $\binom nk$.\qed
\end{corollary}

This proves Theorem~\ref{thm:local}.

\section{The explicit lower bound}\label{sec:lower}

\begin{proposition}\label{prop:5}
Let $n$ be admissible with bad index $i_0$, and let $p^a\|k$. Then $n-i_0$ is the unique element of the block with maximal $p$-adic valuation, and
\[v_p(n-i_0)\ \ge\ a+\lfloor\log_p(k-1)\rfloor\ =\ a+L_p-[\,k=p^{L_p}\,].\]
\end{proposition}
\begin{proof}
Write $L=L_p$. By Corollary~\ref{cor:4}, $\Bad_p(n)\ne\varnothing$, so $\Bad_p(n)=\{i_0\}$, and by Lemma~\ref{lem:3}(2) $i_0$ is the unique index of maximal valuation and every other index has valuation at most $M_p-d_p$.

Let $e_1$ be the largest $e$ with $0\le e\le L$ such that the block contains at least two multiples of $p^e$ ($e=0$ always qualifies); one of them is not $n-i_0$, so the second-largest valuation in the block is at least $e_1$, whence $M_p-d_p\ge e_1$. For every $e\in(e_1,L]$ the block contains exactly one multiple of $p^e$; since $\lfloor k/p^e\rfloor\ge1$, Lemma~\ref{lem:2} forces $\lfloor k/p^e\rfloor=1$ and $\delta_e=0$.

Suppose $k\ne p^L$. Since $p^a\mid k$, the block contains exactly $k/p^a\ge2$ multiples of $p^a$, so $e_1\ge a$. The levels $e\le a$ (where $k\bmod p^e=0$, so $\delta_e=0$) and the levels $e\in(e_1,L]$ are disjoint sets of non-carries, so $d_p\ge a+(L-e_1)$, and $M_p\ge e_1+d_p\ge a+L$.

If $k=p^L$, then $a=L$, $d_p=L$, and the block contains $p\ge2$ multiples of $p^{L-1}$, so the second-largest valuation is at least $L-1$ and $M_p\ge(L-1)+L=2L-1=a+\lfloor\log_p(k-1)\rfloor$.
\end{proof}

\begin{proof}[Proof of Theorem~\ref{thm:lower}]
By Proposition~\ref{prop:5}, $p^{a+\lfloor\log_p(k-1)\rfloor}\mid n-i_0$ for every $p^a\|k$, so $W'_k\mid n-i_0$ and $n_k\ge n_k-i_0\ge W'_k$. The equality cases were checked against Tables~\ref{tab:values} and~\ref{tab:values2} and the OEIS values: $n_4-1=8$, $n_6-3=72$, $n_{14}-9=784$, and $n_k-i_0>W'_k$ for every other $k\le140$.
\end{proof}

\noindent\emph{Formal verification.} Proposition~\ref{prop:5} (in the slightly stronger form that uses only the divisibility of the elements $n-i$, $i\ne i_0$), the divisibility $W'_k\mid n-i_0$, and the unconditional bound $W'_k\le n_k$ for $k\ge3$ have been formally verified in Lean~4 with Mathlib, using the definition of $n_k$ from \cite{FC} verbatim. The same file proves rickyc's sharper divisibility $p^{\,a+v_p\binom{k-1}{e}+\lfloor\log_p\max(e,k-1-e)\rfloor}\mid n-e$ for every $p^a\|k$, and that the failing index is $e=n\bmod k$; together these are the structural input of the search of jbell31, whose minimality claims therefore no longer rest on an unformalized step. It also proves Cambie's bound $n_k\le k\operatorname{lcm}(1,\dots,k-1)$ independently of \cite{Ku26}; it is part of \cite{code}. The case $k=2$ is immediate ($W'_2=2<4=n_2$). The statements $W'_k\le n_k$ ($k\ge2$) and $e=n\bmod k$ with $D(k,e)\mid n-e$ are registered, with their proofs, in the Palomar registry of machine-checked mathematics as \texttt{PALOMAR-2026-10-23-000014} (repository \texttt{Paulee314/erdos1063-lean}), after independent kernel checks and an automated editorial review.

\begin{remark}
(i) Since $\lfloor\log_p(k-1)\rfloor\ge\lfloor\log_p(k/2)\rfloor$, this improves the $k/2$ bound from the discussion thread (as does rickyc's $k-1$ version there), with strict improvement when $p^{L_p}>k/2$ and $k\ne p^{L_p}$: for $k=6,14,18,54$, $W'_k=72,784,2592,46656$ against $36,392,1296,23328$.
(ii) For prime $k$ the bound is only $n_k\ge k$. It is not comparable with the claimed $c(\log k)^2$ of \eqref{eq:cip} without an explicit $c$; the data force $c\le0.740$ (from $k=19$), and with that value $\log W'_k$ exceeds $c(\log k)^2$ for $28$ of the $80$ composite $k\le110$, those with several small prime factors. Its merit is that it is explicit and attained, not that it is large.
(iii) Along primorials $k=p_1\cdots p_m$, $W'_k>k^{\omega(k)}k/\rad(k)$ gives $n_k\ge\exp\big((1+o(1))(\log k)^2/\log\log k\big)$.
\end{remark}

\section{Primes larger than \texorpdfstring{$\sqrt k$}{sqrt k}}\label{sec:large}

\begin{lemma}\label{lem:6}
Let $p$ be a prime with $\sqrt k<p\le k$, so $L_p=1$ and $p^2>k$, and let $j=\lfloor k/p\rfloor$.
\begin{enumerate}
\item If $n\bmod p<k\bmod p$, then $\Bad_p(n)=\varnothing$.
\item Otherwise, if $j=1$, the block contains exactly one multiple of $p$, at index $n\bmod p$, and $\Bad_p(n)=\{n\bmod p\}$.
\item Otherwise ($j\ge2$), the block contains exactly $j$ multiples of $p$, at most one of which is divisible by $p^2$; $|\Bad_p(n)|=1$ iff one of them is divisible by $p^2$, and then that index is the bad one; if none is, all $j$ are bad.
\end{enumerate}
Consequently \eqref{eq:closed} holds.
\end{lemma}
\begin{proof}
With $L_p=1$, $d_p=[\,n\bmod p\ge k\bmod p\,]$. If $n\bmod p<k\bmod p$ then $d_p=0$ and $\Bad_p=\varnothing$ by Lemma~\ref{lem:3}(1). Otherwise $d_p=1$, so $\Bad_p=\{i: v_p(n-i)\ge M_p\}$ is the set of indices of maximal valuation, and by Lemma~\ref{lem:2} the block contains exactly $j+\delta_1=j$ multiples of $p$. If $j=1$, the unique multiple of $p$ has index $i\equiv n\pmod p$ with $0\le i<p\le k$, i.e.\ $i=n\bmod p$; its valuation is $M_p\ge1$ and all other valuations are $0$. If $j\ge2$, two multiples of $p^2$ in the block would differ by at least $p^2>k$, so at most one index has valuation $\ge2$; if there is one it is the unique maximum, and otherwise all $j$ multiples of $p$ tie at valuation $1=M_p$.

For \eqref{eq:closed}: the events $d_p=0$ and $d_p=1$ are disjoint, and $\Pr(d_p=0)=(k\bmod p)/p$. Given $d_p=1$: if $j=1$, $\Bad_p=\{i_0\}$ iff $n\bmod p=i_0$, which requires $i_0<p$ and $i_0\ge k\bmod p$ and has probability $1/p$ (if $i_0\ge p$ then $i_0\bmod p=i_0-p<k-p=k\bmod p$, so the indicator in \eqref{eq:closed} vanishes, as it should); if $j\ge2$, $\Bad_p=\{i_0\}$ iff $p^2\mid n-i_0$, which has probability $1/p^2$ and implies $n\bmod p=i_0\bmod p$, so that $d_p=1$ is exactly the condition $i_0\bmod p\ge k\bmod p$.
\end{proof}

\section{Periodicity and the exact density}\label{sec:density}

\begin{proof}[Proof of Theorem~\ref{thm:density}]
\emph{Step 1: $\Bad_p(n)$ depends only on $n\bmod p^{2L}$, where $L=L_p$.} Let $n'\equiv n\pmod{p^{2L}}$. The deficiency $d_p$ depends only on $n\bmod p^L$, so $d_p(n')=d_p(n)$. For each index $i$, $n'-i\equiv n-i\pmod{p^{2L}}$, so $v_p(n'-i)=v_p(n-i)$ whenever $v_p(n-i)<2L$, and $v_p(n'-i)\ge2L$ iff $v_p(n-i)\ge2L$. Two indices with valuation $\ge L+1$ would differ by a multiple of $p^{L+1}>k$, so at most one index has valuation $\ge L+1$. If no index has valuation $\ge2L$, all valuations agree between $n$ and $n'$, and so do $M_p$, $d_p$ and $\Bad_p$. If one index $i^*$ has $v_p(n-i^*)\ge2L$, it is the unique maximum for both $n$ and $n'$, and every other valuation is $\le L$ and agrees; by Lemma~\ref{lem:3}, $i^*$ is bad iff $d_p\ge1$, for $n$ and $n'$ alike, while another index $i$ would need $v_p(n-i)>M_p-d_p\ge2L-L=L$, which is impossible. So $\Bad_p(n')=\Bad_p(n)$, and $E_{p,i_0}$ is a union of residue classes modulo $p^{2L}$.

The modulus $p^{2L-1}$ does not suffice in general. For $(k,p)=(6,3)$: the block of $n=12$ is $\{7,\dots,12\}$ with $v_3(9)=2$ at $i=3$ and $v_3(12)=1$ at $i=0$, so $M_3=2$, $d_3=[12\bmod3\ge6\bmod3]=1$ and $\Bad_3=\{3\}$; the block of $n=15\equiv12\pmod3$ is $\{10,\dots,15\}$ with $v_3(12)=v_3(15)=1$, so $M_3=1$, $d_3=1$ and $\Bad_3=\{0,3\}$.

\emph{Step 2: admissibility.} By Lemma~\ref{lem:1} the set of bad indices is $\bigcup_{p\le k}\Bad_p(n)$, and it is nonempty by Corollary~\ref{cor:4}. So $n$ is admissible iff there is an $i_0$ with $\Bad_p(n)\subseteq\{i_0\}$ for every $p\le k$, i.e.\ $n\in\bigcap_{p\le k}E_{p,i_0}$. If this held for two indices $i_0\ne i_1$, every $\Bad_p(n)$ would be empty, contradicting Corollary~\ref{cor:4}; so the $k$ sets $\bigcap_pE_{p,i_0}$ are pairwise disjoint.

\emph{Step 3: density.} By Step 1 and the Chinese remainder theorem, $\bigcap_{p\le k}E_{p,i_0}$ is a union of residue classes modulo $D_k$, and the proportion of classes it contains is $\prod_{p\le k}\Pr(E_{p,i_0})$, since $n\bmod D_k\mapsto(n\bmod p^{2L_p})_{p\le k}$ is a bijection. The natural density of a periodic set is its proportion of classes per period, and by Step 2 the admissible $n$ form the disjoint union over $i_0$; the restriction $n\ge2k$ does not affect density. The closed form \eqref{eq:closed} is Lemma~\ref{lem:6}.
\end{proof}

For $p\le\sqrt k$ the factor $\Pr(E_{p,i_0})$ can be computed by enumerating $n\bmod p^{2L_p}$; our implementation enumerates $n\bmod p^{L_p+1}$ and treats the valuation of the at most one multiple of $p^{L_p+1}$ as a geometric tail. Two independent implementations agree, and brute-force counts over full periods confirm, for instance, $\delta(7)=269/2520$, $\delta(11)=9803/831600$ and $\delta(12)=1/1386$.

\section{The size of the density}\label{sec:asymp}

Throughout this section $p$ denotes a prime, $\{x\}=x-\lfloor x\rfloor$, and for a prime $p\in(\sqrt k,k]$ with $p\nmid k$ we write
\[j=j(p)=\lfloor k/p\rfloor,\qquad r=r(p)=k-jp=k\bmod p\in[1,p-1],\]
so that $\{k/p\}=r/p$ and $1\le j<\sqrt k$. The at most one prime $p>\sqrt k$ dividing $k$ is excluded from $S(k)$ because $k\bmod p=0$; Part~1 below shows it contributes $O(\log k)$ to $\log(1/\delta(k))$. We use three classical inputs: the prime number theorem in the form $\pi(x)=\operatorname{li}(x)+E(x)$ with $E(x)\ll x\exp(-c\sqrt{\log x})$ \cite{dlVP}; the Brun--Titchmarsh inequality in the form of Montgomery and Vaughan \cite{MV73}, $\pi(x+y)-\pi(x)\le2y/\log y$ for integers $x\ge0$, $y\ge2$, which for a real interval of length $y\ge2$ (contained in an integer interval of length at most $y+1$) gives
\begin{equation}\label{eq:BT}
\pi(x+y)-\pi(x)\le\frac{3y}{\log y}\qquad(x\ge0,\ y\ge2\text{ real}),
\end{equation}
since $2(y+1)/\log(y+1)\le3y/\log y$ for $y\ge2$; and Chebyshev's bounds $\pi(x)\ll x/\log x$, $\vartheta(x)\ll x$.

\begin{lemma}[injectivity]\label{lem:A}
The map $p\mapsto k\bmod p$ is injective on the primes $p\in(\sqrt k,k]$ with $p\nmid k$.
\end{lemma}
\begin{proof}
If $p\ne q$ are two such primes with $k\bmod p=k\bmod q=r$, then $pq\mid k-r$; but $k-r=jp\ge p>0$ and $k-r\le k<pq$.
\end{proof}

\begin{lemma}[the constant]\label{lem:B}
For $j\ge1$, $c_j:=\int_j^{j+1}\frac{-\log(x-j)}{x^2}\,dx=\frac1j\log\big(1+\frac1j\big)$, and $\sum_{j\ge1}c_j=C$ with $\sum_{j>J}c_j\le1/J$.
\end{lemma}
\begin{proof}
Substituting $u=x-j$, $c_j=\int_0^1\frac{-\log u}{(j+u)^2}\,du$. Integrate by parts with $v(u)=\frac{u}{j(j+u)}$, which satisfies $v'(u)=(j+u)^{-2}$ and $v(0)=0$:
\[c_j=\Big[-\log u\cdot\frac{u}{j(j+u)}\Big]_0^1+\int_0^1\frac1u\cdot\frac{u}{j(j+u)}\,du=\frac1j\int_0^1\frac{du}{j+u}=\frac1j\log\frac{j+1}{j},\]
the boundary term vanishing since $\log1=0$ and $u\log u\to0$. As $\log(1+1/j)\le1/j$, $c_j\le j^{-2}$ and $\sum_{j>J}c_j\le1/J$; splitting $[1,\infty)$ at the integers gives $\sum_jc_j=\int_1^\infty-\log\{x\}\,x^{-2}dx$. The expansion $\log(1+1/j)=\sum_{m\ge1}(-1)^{m+1}j^{-m}/m$ gives the zeta form of $C$: the $j=1$ term is $\log2=\sum_m(-1)^{m+1}/m$ (conditionally convergent), and for $j\ge2$ the double series converges absolutely, since $\zeta(m+1)-1\le3\cdot2^{-m-1}$, so the sums may be interchanged.
\end{proof}

\subsection*{Part 1: \texorpdfstring{$\log(1/\delta(k))=S(k)+O(\sqrt k)$}{log(1/delta) = S(k) + O(sqrt k)}}

\emph{Upper bound for $\delta$.} In \eqref{eq:delta}, bound each of the $k$ summands by the product over the primes $p>\sqrt k$, $p\nmid k$ only (the other factors are at most $1$), and each such factor, using \eqref{eq:closed}, by $\{k/p\}+1/p$:
\[\delta(k)\le k\prod_{\substack{\sqrt k<p\le k\\p\nmid k}}\Big(\Big\{\frac kp\Big\}+\frac1p\Big)=k\prod\Big\{\frac kp\Big\}\prod\Big(1+\frac1{k\bmod p}\Big).\]
By Lemma~\ref{lem:A} the values $k\bmod p$ are distinct positive integers at most $k$, so $\sum\log(1+\frac1{k\bmod p})\le\sum_{r\le k}\frac1r\le\log k+1$. Hence $\log(1/\delta(k))\ge S(k)-2\log k-1$.

\emph{Lower bound for $\delta$.} Keep only the term $i_0=0$ in \eqref{eq:delta}. For $p\le\sqrt k$ (so $L=L_p\ge1$): if $p^{2L}\mid n$ then $v_p(n)\ge2L$ while $v_p(n-i)=v_p(i)\le L<2L$ for $1\le i<k$ (as $i<k<p^{L+1}$); so $M_p=v_p(n)$ is attained only at $i=0$, and every other index has valuation $\le L\le M_p-d_p$ because $d_p\le L$; by Lemma~\ref{lem:3}, $\Bad_p(n)\subseteq\{0\}$. Thus $\Pr(E_{p,0})\ge p^{-2L}\ge k^{-2}$, and $\sum_{p\le\sqrt k}2\log k\ll\sqrt k$ by Chebyshev. For $\sqrt k<p\le k$ with $p\nmid k$, \eqref{eq:closed} gives $\Pr(E_{p,0})\ge\{k/p\}$. For the at most one prime $p>\sqrt k$ with $p\mid k$, $k\bmod p=0$ and \eqref{eq:closed} gives $\Pr(E_{p,0})\ge1/p^2\ge1/k^2$. Hence $\log(1/\delta(k))\le S(k)+O(\sqrt k)$.

\subsection*{Part 2: \texorpdfstring{$S(k)=(C+o(1))\,k/\log k$}{S(k) = (C+o(1)) k/log k}}

Fix $\varepsilon\in(0,1/e)$ and an integer $J\ge2$; they stay fixed while $k\to\infty$ (assume $k>(J+1)^2$), and are sent to $0$ and $\infty$ at the end. Cover the primes $p\in(\sqrt k,k]$, $p\nmid k$, by three classes according to $r=k\bmod p$ and $j=\lfloor k/p\rfloor$:
\[\mathcal T:\ r\le4j,\qquad\mathcal D:\ 4j<r\le\varepsilon p,\qquad\mathcal B:\ r>\varepsilon p.\]
($\mathcal T$ and $\mathcal B$ can overlap for $p<2\sqrt{k/\varepsilon}$. This is harmless: the upper bound on $S(k)$ only needs a cover, and the lower bound uses only the part of $\mathcal B$ with $j\le J$, which is a subset of the terms.)

\emph{Class $\mathcal T$.} For primes with $r\le4j$ we have $jp\le k-1$ (as $r\ge1$) and $k\le jp+4j$, so $p\in[k/j-4,\,k/j)$, a half-open interval of length $4$ containing at most $4$ integers. So there are at most $4$ such primes for each of the fewer than $\sqrt k$ values of $j$, and every term is at most $\log p\le\log k$. The class contributes at most $4\sqrt k\log k=o(k/\log k)$.

\emph{Class $\mathcal D$.} Group the primes by $j$ and by dyadic ranges $r\in(R,2R]$ with $R$ a power of $2$. A prime in the group satisfies $k-2R\le jp<k-R$, so $p$ lies in an interval of length $R/j$; and $r>4j$ forces $R>2j$, so $R/j>2$ and \eqref{eq:BT} applies: the group has at most $\frac{3R/j}{\log(R/j)}$ primes. Each term is $\log(p/r)\le\log\frac{k}{jR}$, since $p\le k/j$ and $r>R$. Also $R<r\le\varepsilon p\le\varepsilon k/j$.

Put $V=\varepsilon k/j^2$ and, for a group with range $(R,2R]$, $u=R/j$; then $2<u<V$ and $\log\frac{k}{jR}=\log\frac Vu+\log\frac1\varepsilon$. So the contribution of fixed $j$ is at most
\[\sum_u\frac{3u}{\log u}\Big(\log\frac Vu+\log\frac1\varepsilon\Big),\]
where $u$ runs over the values $R/j$ with $R$ a power of two and $2<u<V$. The admissible $R$ are the powers of two below $\varepsilon k/j$; the largest, $R_{\max}$, satisfies $R_{\max}\ge\varepsilon k/(2j)$, so the $u$ are $V'/2^m$ ($m\ge0$) with $V/2\le V'<V$, and $V/2^{m+1}\le u\le V/2^m$.
\begin{itemize}
\item For $m\le\frac12\log_2V-1$: here $\log u\ge\log(V/2^{m+1})\ge\frac12\log V$ and $\log(V/u)\le(m+1)\log2$, so these terms sum to at most $\frac{6V}{\log V}\sum_{m\ge0}\frac{(m+1)\log2+\log(1/\varepsilon)}{2^m}\le\frac{24V}{\log V}\big(\log2+\log\frac1\varepsilon\big)$.
\item For $m>\frac12\log_2V-1$: here $u\le V/2^m<2\sqrt V$, while $u>2$ gives $\log u\ge\log2$ and $\log(V/u)\le\log V$. Each such group contributes at most $\frac{6\sqrt V}{\log2}\big(\log V+\log\frac1\varepsilon\big)$, and there are at most $\log_2V$ such groups (as $u>2$ forces $2^m<V/2$). So these terms sum to $\ll\sqrt V(\log k)^2$, as $V\le k$ and $\varepsilon$ is fixed.
\end{itemize}
Now sum over $j$; only $j$ with $V>2$, i.e.\ $j<\sqrt{\varepsilon k/2}$, contribute. The $\sqrt V$ terms give $\ll\sqrt{\varepsilon k}(\log k)^2\sum_{j<\sqrt k}\frac1j\ll\sqrt k(\log k)^3=o(k/\log k)$. For $j\le k^{1/4}$, $\log V\ge\log(\varepsilon\sqrt k)\ge\frac14\log k$ for large $k$, so $\sum_{j\le k^{1/4}}\frac{24V}{\log V}(\log2+\log\frac1\varepsilon)\le\frac{96\varepsilon k}{\log k}(\log2+\log\frac1\varepsilon)\sum_jj^{-2}\ll\varepsilon\log(1/\varepsilon)\frac{k}{\log k}$. For $k^{1/4}<j<\sqrt{\varepsilon k/2}$, $\log V\ge\log2$, so these give $\ll\varepsilon k\log(1/\varepsilon)\sum_{j>k^{1/4}}j^{-2}\ll\varepsilon\log(1/\varepsilon)k^{3/4}=o(k/\log k)$. Altogether, class $\mathcal D$ contributes at most $A\,\varepsilon\log(1/\varepsilon)\,k/\log k+o(k/\log k)$ for an absolute constant $A$.

\emph{Class $\mathcal B$, $j>J$.} Each term is $\log(p/r)<\log(1/\varepsilon)$. The primes with $j>J$ satisfy $p<k/J$, so by \eqref{eq:BT} with $x=0$ there are at most $\pi(k/J)\le\frac{3k}{J\log(k/J)}\le\frac{6k}{J\log k}$ of them for $k\ge J^2$. The contribution is at most $6\log(1/\varepsilon)\,k/(J\log k)$.

\emph{Class $\mathcal B$, $j\le J$: the main term.} Fix $j\le J$. The primes of class $\mathcal B$ with $\lfloor k/p\rfloor=j$ are the primes in the interval
\[I_j(\varepsilon)=\Big(\frac k{j+1},\ \frac k{j+\varepsilon}\Big),\]
because $r>\varepsilon p\iff k>(j+\varepsilon)p$. (A prime $p\mid k$ with $\lfloor k/p\rfloor=j$ equals $k/j$, which is not in $I_j(\varepsilon)$, so the condition $p\nmid k$ excludes nothing here; and $I_j(\varepsilon)\subset(\sqrt k,k]$ since $k>(J+1)^2$.) On this interval put $g(t)=\log\frac{t}{k-jt}=-\log\big(\frac kt-j\big)$; then $0<g(t)<\log(1/\varepsilon)$ and $g'(t)=\frac1t+\frac j{k-jt}\le\frac{1+j/\varepsilon}t\le\frac{(J+1)(1+J/\varepsilon)}k$.

By partial summation with $\pi(t)=\operatorname{li}(t)+E(t)$, writing $a<b$ for the endpoints of $I_j(\varepsilon)$,
\[\sum_{p\in I_j(\varepsilon)}g(p)=\int_a^b\frac{g(t)}{\log t}\,dt+\big[g(t)E(t)\big]_a^b-\int_a^bg'(t)E(t)\,dt+O\big(\log\tfrac1\varepsilon\big),\]
where the $O(\cdot)$ covers the endpoints. The two error terms are at most $2\log(1/\varepsilon)\max_{[a,b]}|E|+(b-a)\max_{[a,b]}|g'|\max_{[a,b]}|E|\ll_{\varepsilon,J}k\exp(-c'\sqrt{\log k})=o(k/\log k)$, since $[a,b]\subset[k/(J+1),k]$. For the main integral, $\log t=\log k+O(\log(J+1))=(1+o(1))\log k$ on $[a,b]$, so, substituting $x=k/t$, $dt=-k\,x^{-2}dx$,
\[\int_a^b\frac{g(t)}{\log t}\,dt=\frac{(1+o(1))\,k}{\log k}\int_{j+\varepsilon}^{j+1}\frac{-\log(x-j)}{x^2}\,dx=:\frac{(1+o(1))\,k}{\log k}\,c_j(\varepsilon).\]
Here $0\le c_j-c_j(\varepsilon)=\int_j^{j+\varepsilon}\frac{-\log(x-j)}{x^2}dx\le\frac1{j^2}\int_0^\varepsilon-\log u\,du=\frac{\varepsilon(1+\log(1/\varepsilon))}{j^2}\le\frac{2\varepsilon\log(1/\varepsilon)}{j^2}$ for $\varepsilon<1/e$. Summing over $j\le J$, with Lemma~\ref{lem:B} and $\sum_jj^{-2}<2$,
\[\sum_{\substack{p\in\mathcal B\\j(p)\le J}}\log\frac pr=\frac{(1+o(1))\,k}{\log k}\sum_{j\le J}c_j(\varepsilon),\qquad C-\frac1J-4\varepsilon\log\frac1\varepsilon\ \le\ \sum_{j\le J}c_j(\varepsilon)\ \le\ C.\]

\emph{Conclusion.} All terms of $S(k)$ are positive, so the $j\le J$ part of class $\mathcal B$ alone gives the lower bound, and the four pieces together give the upper bound:
\[\Big(C-\frac1J-4\varepsilon\log\frac1\varepsilon-o(1)\Big)\frac k{\log k}\ \le\ S(k)\ \le\ \Big(C+A\varepsilon\log\frac1\varepsilon+\frac{6\log(1/\varepsilon)}J+o(1)\Big)\frac k{\log k},\]
where $o(1)\to0$ as $k\to\infty$ with $\varepsilon,J$ fixed. Letting $k\to\infty$, then $J\to\infty$, then $\varepsilon\to0$, gives $S(k)=(C+o(1))\,k/\log k$. With Part~1 this proves Theorem~\ref{thm:asymp}.\qed

\begin{remark}
The convergence is slow. Averaged over $21$ consecutive values of $k$, $S(k)/(k/\log k)$ is $1.16$ near $k=100$, $1.37$ near $10^4$, and $1.34$ near $10^6$; the next term should be $c_2k/\log^2k$ with $c_2=\int_1^\infty-\log\{x\}\log x\,x^{-2}dx=1.0016551589\ldots$, obtainable by keeping the variation of $\log t$ in the main integral. The class bounds are visible numerically: the class $\mathcal D$ sum divided by $\varepsilon\log(1/\varepsilon)k/\log k$ is $2.29,2.08,2.10,1.97$ at $(k,\varepsilon)=(10^6,0.1),(10^6,0.02),(10^7,0.02),(10^7,0.005)$, so $A\approx2$ in practice. The small-prime part $\log(1/\delta(k))-S(k)$ lies between $3.2$ and $17.0$ for $31\le k\le110$.
\end{remark}

\section{The growth conjecture and the evidence}\label{sec:evidence}

Theorem~\ref{thm:asymp} concerns the density of admissible $n$. If the admissible set behaved like a random set of density $\delta(k)$, its least element would be about $\Exp(1)/\delta(k)$; this is Conjecture~\ref{conj}. It is not a theorem, and we do not know how to prove either inequality in it (see Section~\ref{sec:cip}). The evidence is as follows.

\emph{New values.} Table~\ref{tab:values} gives $n_k$ for $60\le k\le110$, and Table~\ref{tab:values2} for $111\le k\le140$ and the primes $149\le k\le167$, extending the OEIS table. The values for $60\le k\le75$ were posted earlier in the discussion thread by jbell31 (28 July 2026), from a search that assumes rickyc's divisibility result (now formally verified; see the end of Section~\ref{sec:lower}); ours were obtained independently and agree with all $16$. The values for $k\le110$ were obtained by an exhaustive search over the candidate set defined by the necessary conditions of Proposition~\ref{prop:5} and Lemma~\ref{lem:6}, with every candidate tested by the exact criterion of Lemmas~\ref{lem:1}--\ref{lem:2}; the same program reproduces the $58$ OEIS values for $k\le59$, and minimality was confirmed by a second, independently written exhaustive search, and for several $k$ by a plain scan from $2k$. The values for $k\ge111$ come from a faster exhaustive search. For each possible failing index $e$ it scans only $n\equiv e\pmod{D(k,e)}$, where $D(k,e)=\prod_{p^a\|k}p^{\,a+v_p\binom{k-1}{e}+\lfloor\log_p\max(e,k-1-e)\rfloor}$ is rickyc's sharper modulus (this divisibility and $e=n\bmod k$ are formally verified; see the end of Section~\ref{sec:lower}), and it applies Lemma~\ref{lem:6} at every prime in $(\sqrt k,k]$; again every candidate gets the exact test. Before the new runs it reproduced all $109$ values for $2\le k\le110$, and the larger tables used for $k\ge121$ were re-checked on $14$ known values. The second, independently written search confirmed minimality for $111\le k\le119$ and $k=127$; for the other new values it is too slow, and their minimality rests on the single search. Every value was verified with exact integer arithmetic (exactly one of $n,\dots,n-k+1$ fails to divide $\binom nk$).

\emph{Prime versus composite $k$.} The model reproduces the large swings between neighbouring $k$: primes dividing $k$ force $n-i_0$ to be divisible by $p^{a+\lfloor\log_p(k-1)\rfloor}$ (Proposition~\ref{prop:5}), which is why $n_{30}\approx3.8\cdot10^7$ while $n_{31}=8100$.

\emph{Fit.} Let $x_k=\log n_k-\log(1/\delta(k))$. For $\log$ of an $\Exp(1)$ variable the mean is $-\gamma=-0.577$ and the standard deviation $\pi/\sqrt6=1.283$. For $31\le k\le140$ ($110$ values; sample standard deviations throughout) $x_k$ has mean $-0.59$ and standard deviation $1.20$, a Kolmogorov--Smirnov statistic $D=0.055$ against $\log\Exp(1)$ ($5\%$ critical value $0.129$), and $\log n_k$ and $\log(1/\delta(k))$ have correlation $0.993$ with regression slope $0.994$; for $31\le k\le110$ alone the figures were $-0.58$, $1.19$, $D=0.083$ and $0.988$. Adding the five primes $149\le k\le167$ changes each figure by less than $0.02$. For $k\le30$ the fit is poor (mean $+0.19$, $D=0.348$); there the constraint $n\ge2k$ and the coarse structure of the admissible set matter.

\emph{Out-of-sample tests.} Before computing any $n_k$ with $k\ge60$, the central $90\%$ interval of $\Exp(1)/\delta(k)$, namely $[0.0513/\delta(k),\,2.996/\delta(k)]$, was recorded for every $60\le k\le100$; $38$ of the $41$ values fell inside (outside: $k=64,91,95$). The same was then done for $101\le k\le110$, with $9$ of $10$ inside (outside: $k=105$). Combined: $47$ of $51$, against $45.9$ expected. Two further tests were registered in the same way, each before its values were computed. For $111\le k\le120$ and the primes $113,127,131,137,139$, all $14$ values fell inside ($12.6$ expected). For the remaining $121\le k\le140$ and the primes $149\le k\le199$, the $21$ values computed so far ($k\le140$ and $149\le k\le167$) gave $17$ inside ($18.9$ expected; outside: $k=123,132$ below and $k=140,151$ above). For $k=173$ the search was stopped after checking every $n<2.21\cdot10^{22}$ without finding an admissible $n$, so $n_{173}>2.2\cdot10^{22}$, which is consistent with that prediction's interval $[1.05\cdot10^{21},\,6.15\cdot10^{22}]$; the larger primes were not run. Over all four tests: $78$ of $86$, against $77.4$ expected. The model has no fitted parameters, so the order of events could not have influenced the predictions; the timestamped prediction files are in \cite{code}.

\emph{Beyond the least element.} For $k\in\{40,45,50,61,64,67\}$ we enumerated the first $139$--$168$ admissible $n$ (Table~\ref{tab:traj}). The $m$-th admissible number is close to $m/\delta(k)$ (median of $m\,\delta(k)\,n^{(m)}$ between $0.90$ and $1.22$), the count below $X$ is within $11\%$ of $\delta(k)X$, the gaps between consecutive failing elements $n-i_0$ have coefficient of variation $0.84$--$1.02$ (exponential: $1$), and the observed frequencies of the failing index $i_0$ match the model's predicted weights $\prod_p\Pr(E_{p,i_0})$. So the model describes the whole admissible set, not only its least element; the least element is the noisy statistic (for $k=64$ it is $20$ times below $1/\delta$, the outlier $x_{64}=-3.0$ above). For $k\le30$ the admissible $n$ cluster (e.g.\ $n_{30}=37584001$ and $37584002$ share the failing element $37584000$), which is consistent with the worse fit at small $k$; for $k\ge40$ clusters are rare. Finally, an admissible $n$ lands in the carry-free window $n\bmod p<k\bmod p$ at all but on average $0.4$--$0.7$ of the primes $p\in(\sqrt k,k]$, $p\nmid k$, the exceptions being almost always primes $p>k/2$ dividing $n-i_0$. This is the structure a construction would have to reproduce (Section~\ref{sec:cip}).

\begin{table}[h]
\caption{The first admissible $n$ beyond $n_k$: $N$ values enumerated, the median of $m\,\delta(k)\,n^{(m)}$ over $m\le N$ (model: $\approx1$), the ratio of the count below $X=n^{(N)}$ to $\delta(k)X$, the coefficient of variation of the gaps between consecutive failing elements, and the mean number of large primes ($\sqrt k<p\le k$, $p\nmid k$) at which $n$ is \emph{not} in the carry-free window.}\label{tab:traj}
{\setlength{\tabcolsep}{8pt}\begin{tabular}{rrrrrr}
\toprule
$k$ & $N$ & median $m\delta n^{(m)}$ & count$/\delta X$ & gap CV & non-window primes per $n$ \\
\midrule
40 & 139 & 1.05 & 0.95 & 0.97 & 0.43 \\
45 & 139 & 1.22 & 0.94 & 0.90 & 0.67 \\
50 & 150 & 1.08 & 0.98 & 0.94 & 0.48 \\
61 & 168 & 0.90 & 1.11 & 0.94 & 0.66 \\
64 & 150 & 0.99 & 1.00 & 1.02 & 0.54 \\
67 & 161 & 0.98 & 1.06 & 0.84 & 0.41 \\
\bottomrule
\end{tabular}}
\end{table}


\section{Relation to the claimed bounds of Cipollini}\label{sec:cip}

If both claims in \eqref{eq:cip} hold, then Conjecture~\ref{conj} sits strictly between them and says that the upper bound is off by a factor of order $\log\log k$. In the data, for $31\le k\le140$, $\log n_k$ is between $0.47$ and $0.89$ (mean $0.66$) of the main term $\frac k{\log k}(\log\log k+\log\log\log k+\log2)$.

The local lemmas of \cite{Ci26a} (its Lemma~2, ``stabilization'', and Lemma~3, ``one carry'') are the same facts as our Lemma~\ref{lem:3} and Lemma~\ref{lem:6}; Lemma~3 of \cite{Ci26b}, that $k\mid n-i_0$ for admissible $n$, is a special case of Proposition~\ref{prop:5}. The construction in \cite{Ci26a} takes $n=e+Mt$ with a small exceptional index $e\approx k/(2\log k\log\log k)$, forces the primes dividing $k\binom{k-1}e$ by divisibility, and uses a pigeonhole (simultaneous approximation) to put $n\bmod p$ into the carry-free window $[0,k\bmod p)$ for the remaining large primes. Pigeonhole on differences yields symmetric windows, of half-width $\min(e,(k\bmod p)-e)\le e$, at a cost of $\log(p/e)$ per prime; summed over primes this is $\frac k{\log k}(\log(k/e)-1+o(1))$. The density model's natural window is the full $[0,k\bmod p)$, at cost $\log\frac{p}{k\bmod p}$ per prime, and these costs sum to $S(k)\sim Ck/\log k$. The enumeration in Section~\ref{sec:evidence} shows that actual admissible $n$ do sit in the one-sided window at essentially every large prime. The gap between the two costs is exactly the gap between what a symmetric pigeonhole can reach and what a one-sided target would require; closing it appears to need an equidistribution-type input for the vectors $(n\bmod p)_{p}$ over $n\le e^{(C+\varepsilon)k/\log k}$, a range far smaller than the product of the moduli, and we do not know how to obtain one. The lower-bound half of the conjecture, that no admissible $n$ lies below $e^{(C-\varepsilon)k/\log k}$, is a least-element problem of Linnik type and seems out of reach.

\appendix
\section{The values \texorpdfstring{$n_k$}{n\_k} for \texorpdfstring{$60\le k\le140$}{60 <= k <= 140} and prime \texorpdfstring{$k\le167$}{k <= 167}}

\begin{table}[h]
\caption{$n_k$ and the bad index $i_0$ (so $n_k-i_0$ is the element of the block that fails to divide $\binom{n_k}{k}$).}\label{tab:values}
\footnotesize
\begin{tabular}{rrr@{\quad}rrr@{\quad}rrr}
\toprule
$k$ & $n_k$ & $i_0$ & $k$ & $n_k$ & $i_0$ & $k$ & $n_k$ & $i_0$\\
\midrule
60 & 2\,117\,441\,088\,029 & 29 & 61 & 569\,181\,900 & 50 & 62 & 155\,275\,274\,056 & 8 \\
63 & 216\,949\,362\,182 & 35 & 64 & 2\,600\,001\,558 & 22 & 65 & 1\,936\,512\,990\,784 & 34 \\
66 & 1\,618\,958\,347\,825 & 49 & 67 & 1\,178\,168\,706 & 37 & 68 & 260\,686\,545\,186 & 34 \\
69 & 709\,479\,687\,556 & 28 & 70 & 66\,383\,308\,992\,066 & 66 & 71 & 13\,760\,436\,550 & 35 \\
72 & 890\,473\,430\,071 & 55 & 73 & 40\,861\,892\,736 & 21 & 74 & 3\,146\,202\,243\,081 & 9 \\
75 & 32\,639\,932\,074\,432 & 57 & 76 & 16\,233\,110\,681\,625 & 25 & 77 & 18\,874\,715\,290\,176 & 53 \\
78 & 780\,422\,752\,841\,501 & 29 & 79 & 420\,957\,410\,952 & 14 & 80 & 60\,137\,833\,600\,000 & 0 \\
81 & 2\,418\,072\,694\,480 & 46 & 82 & 94\,878\,375\,215\,625 & 9 & 83 & 885\,437\,574\,400 & 11 \\
84 & 3\,765\,171\,612\,545\,280 & 0 & 85 & 1\,718\,626\,725\,667\,152 & 27 & 86 & 129\,713\,628\,198\,225 & 81 \\
87 & 2\,221\,164\,774\,687\,826 & 82 & 88 & 754\,618\,137\,360\,426 & 42 & 89 & 3\,208\,190\,120\,080 & 80 \\
90 & 57\,753\,167\,368\,032\,004 & 4 & 91 & 10\,789\,171\,853\,293\,890 & 56 & 92 & 2\,643\,061\,175\,354\,691 & 67 \\
93 & 7\,201\,679\,510\,714\,780 & 29 & 94 & 2\,544\,131\,371\,430\,937 & 25 & 95 & 213\,706\,678\,047\,814 & 64 \\
96 & 11\,795\,043\,698\,546\,775 & 87 & 97 & 10\,432\,185\,034\,560 & 35 & 98 & 23\,154\,863\,741\,528\,961 & 1 \\
99 & 7\,525\,993\,870\,612\,546 & 43 & 100 & 24\,275\,936\,775\,680\,073 & 73 & 101 & 83\,339\,504\,740\,900 & 4 \\
102 & 455\,022\,751\,908\,409\,400 & 56 & 103 & 80\,855\,126\,198\,850 & 57 & 104 & 28\,292\,410\,857\,118\,753 & 33 \\
105 & 114\,184\,320\,231\,814\,976 & 101 & 106 & 505\,139\,783\,269\,972\,305 & 81 & 107 & 314\,604\,394\,073\,026 & 96 \\
108 & 42\,039\,517\,780\,512\,072 & 72 & 109 & 92\,095\,914\,818\,376 & 95 & 110 & 4\,499\,581\,121\,078\,560\,101 & 101 \\
\bottomrule
\end{tabular}
\end{table}

\begin{table}[h]
\caption{$n_k$ and the bad index $i_0$ for $111\le k\le140$ and the primes $149\le k\le167$ (computed 23 October 2026).}\label{tab:values2}
\footnotesize
\begin{tabular}{rrr@{\qquad}rrr}
\toprule
$k$ & $n_k$ & $i_0$ & $k$ & $n_k$ & $i_0$\\
\midrule
111 & 624\,203\,258\,168\,949\,800 & 5 & 112 & 235\,084\,258\,952\,159\,310 & 78 \\
113 & 1\,154\,291\,595\,269\,232 & 25 & 114 & 3\,058\,893\,621\,505\,762\,601 & 41 \\
115 & 2\,571\,800\,553\,180\,159\,954 & 79 & 116 & 4\,581\,489\,291\,604\,040\,805 & 101 \\
117 & 2\,625\,464\,741\,944\,934\,180 & 83 & 118 & 12\,488\,331\,360\,383\,631\,792 & 48 \\
119 & 1\,075\,680\,624\,503\,443\,492 & 65 & 120 & 349\,159\,236\,017\,505\,408\,068 & 68 \\
121 & 36\,031\,544\,587\,629\,360 & 99 & 122 & 10\,690\,191\,371\,065\,479\,196 & 28 \\
123 & 539\,955\,158\,318\,483\,546 & 56 & 124 & 2\,319\,409\,815\,118\,477\,347 & 35 \\
125 & 165\,058\,245\,202\,193\,772 & 22 & 126 & 1\,612\,403\,449\,350\,274\,216\,000 & 64 \\
127 & 9\,800\,282\,003\,526\,126 & 48 & 128 & 2\,100\,509\,814\,172\,237\,875 & 51 \\
129 & 366\,918\,069\,874\,600\,035\,072 & 72 & 130 & 2\,672\,890\,339\,561\,621\,920\,000 & 0 \\
131 & 1\,225\,479\,431\,226\,053\,250 & 60 & 132 & 123\,904\,568\,785\,495\,246\,880 & 32 \\
133 & 883\,499\,074\,332\,177\,600\,000 & 53 & 134 & 4\,624\,346\,743\,350\,798\,696\,001 & 65 \\
135 & 894\,259\,177\,676\,725\,050\,114 & 114 & 136 & 266\,331\,343\,761\,076\,882\,500 & 68 \\
137 & 3\,995\,236\,738\,822\,581\,000 & 60 & 138 & 48\,518\,376\,432\,039\,910\,149\,250 & 130 \\
139 & 7\,179\,634\,687\,138\,597\,002 & 111 & 140 & 293\,365\,786\,466\,311\,148\,800\,071 & 71 \\
149 & 59\,804\,718\,543\,532\,815\,636 & 14 & 151 & 316\,542\,517\,930\,011\,220\,500 & 139 \\
157 & 48\,048\,676\,933\,717\,376\,280 & 75 & 163 & 377\,518\,312\,793\,973\,563\,136 & 68 \\
167 & 729\,078\,224\,618\,436\,330\,880 & 8 & & & \\
\bottomrule
\end{tabular}
\end{table}

\begin{thebibliography}{99}
\bibitem{Ci26a} R.~Cipollini (with GPT-5.6 Sol), \emph{An upper bound for Erd\H{o}s--Selfridge problem 1063}, proof claim posted at \url{https://www.erdosproblems.com/1063}, submitted 2026-07-27.
\bibitem{Ci26b} R.~Cipollini (with GPT-5.6 Sol), \emph{A lower bound for Erd\H{o}s--Selfridge problem 1063}, proof claim posted at \url{https://www.erdosproblems.com/1063}, submitted 2026-08-04.
\bibitem{code} P.~Lee, \emph{Code, data and prediction files for Erd\H{o}s problem 1063}, \url{https://github.com/Paulee314/erdos1063}, 2026.
\bibitem{dlVP} C.-J.~de la Vall\'ee Poussin, \emph{Sur la fonction $\zeta(s)$ de Riemann et le nombre des nombres premiers inf\'erieurs \`a une limite donn\'ee}, M\'em.\ Couronn\'es Acad.\ Roy.\ Sci.\ Bruxelles \textbf{59} (1899), 1--74.
\bibitem{EP1063} T.~Bloom (ed.), \emph{Erd\H{o}s problem 1063}, \url{https://www.erdosproblems.com/1063}.
\bibitem{ErSe83} P.~Erd\H{o}s and J.~L.~Selfridge, \emph{Problem 6447}, Amer.\ Math.\ Monthly \textbf{90} (1983), 710.
\bibitem{FC} Google DeepMind, \emph{Formal Conjectures}, \texttt{FormalConjectures/ErdosProblems/1063.lean}, \url{https://github.com/google-deepmind/formal-conjectures}.
\bibitem{Gu04} R.~K.~Guy, \emph{Unsolved problems in number theory}, 3rd ed., Springer, 2004, Problem B31.
\bibitem{Ku26} P.-N.~Kung, L.~Song, D.~Hwang, J.~Yoon, C.-L.~Li, S.~Severini, M.~Ol\v{s}\'ak, E.~Lockhart, Q.~V.~Le, B.~Gokturk, T.~Luong, T.~Pfister and N.~Peng, \emph{LEAP: Supercharging LLMs for formal mathematics with agentic frameworks}, arXiv:2606.03303, 2026.
\bibitem{Mo85} J.-M.~Monier, \emph{Solutions of Advanced Problems: 6447}, Amer.\ Math.\ Monthly \textbf{92} (1985), 435--436.
\bibitem{MV73} H.~L.~Montgomery and R.~C.~Vaughan, \emph{The large sieve}, Mathematika \textbf{20} (1973), 119--134.
\bibitem{OEIS} OEIS Foundation, \emph{Sequence A389360}, \url{https://oeis.org/A389360}.
\bibitem{Wh26} P.~White (with Codex), \emph{Erd\H{o}s \#1063: a real upper-bound jump on $n_k$}, \url{https://erdosproblemaday.com/day/1063-one-large-prime}, 2026-07-25; code at \url{https://github.com/pw/erdos1063-upper-bound}.
\end{thebibliography}

\end{document}
